\subsection{Main Results}
\label{sec:main_results}

% Table~\ref{tab:dcd_main} is declared in 4_dcd_datasets.tex for float placement.
Table~\ref{tab:dcd_main} reports the macro results on the 12 evaluation
datasets (8 Unseen and 4 Novel).  DCD obtains 58.76 Dice, 45.17 IoU, 66.02
clDice, and 63.13 HD95, ranking first on all four metrics.  The closest
baselines are SACM and SAM-OCTA, both re-run under our protocol.  DCD
improves over SACM
(56.24/42.88/62.26/69.20) by 2.52 Dice, 2.29 IoU, and 3.76 clDice, and
reduces HD95 by 6.07 pixels; SAM-OCTA is slightly behind SACM on Dice but
stronger on IoU and clDice, and DCD still improves over it by 2.60 Dice,
1.44 IoU, and 2.56 clDice while reducing HD95 by 14.58 pixels.  The strongest SAM-based methods occupy the top ranks, while
the CNN and transformer baselines are substantially lower under the 3-shot
protocol.  DCD also has the smallest seed variance on Dice, IoU, and HD95
among the SAM-based methods, and its clDice variance is tied for the smallest.
These are macro trends: Section~\ref{sec:per_dataset} shows that the gains
are not uniform across datasets.


\begin{table}[!t]
\centering
\caption{Deployment-rule and training-component ablation on the 12
evaluation datasets.  Four-seed absolute means; Dice/IoU/clDice in \%, HD95
in pixels (lower is better).  ``gating + fusion'' is test-time selection
and fusion on the full model's checkpoint; the full model deploys a
validation-selected fixed candidate (slot 3 in our benchmark).  \emph{W/o fusion} has lower HD95, and
\emph{w/o topology} is slightly higher on Dice/IoU.}
\label{tab:dcd_ablation}
\scriptsize
\setlength{\tabcolsep}{2.5pt}
\begin{tabular*}{\columnwidth}{@{\extracolsep{\fill}}lcccc@{}}
\toprule
Setting & Dice $\uparrow$ & IoU $\uparrow$ & clDice $\uparrow$ & HD95 $\downarrow$ \\
\midrule
\multicolumn{5}{l}{\emph{Alternative deployment}}\\
gating + fusion (test time) & 57.54 & 44.06 & 64.96 & 63.10 \\
\midrule
\multicolumn{5}{l}{\emph{Training ablations}}\\
w/o fusion loss & 58.64 & 45.06 & 65.76 & 62.77 \\
native four heads & 57.54 & 44.02 & 64.72 & 64.56 \\
w/o topology supervision & 58.84 & 45.44 & 65.76 & 63.15 \\
\midrule
\textcolor{red!80!black}{\textbf{Full model}} & \textcolor{red!80!black}{\textbf{58.76}} & \textcolor{red!80!black}{\textbf{45.17}} & \textcolor{red!80!black}{\textbf{66.02}} & \textcolor{red!80!black}{\textbf{63.13}} \\
\bottomrule
\end{tabular*}
\end{table}

\subsection{Ablation Study}
\label{sec:ablation}



Table~\ref{tab:dcd_ablation} isolates the deployment rule and the training
components.  The first row uses gating and fusion at test time with the
same trained weights as the full model; replacing this selector with the
fixed slot improves Dice, IoU, and clDice by 1.22, 1.11, and 1.06 points,
while HD95 changes by only 0.03 pixels.  Because the candidate experts are identical,
this comparison directly supports the deployment-bottleneck hypothesis.  The
native four-head control removes role losses and fusion; compared with the
no-fusion variant, which keeps the role losses, the role objectives add 1.10
Dice, 1.04 IoU, and 1.04 clDice, and reduce HD95 by 1.79 pixels.  Because all
heads share the decoder feature, these objectives also regularize the
representation used by the deployed head.  Topology
supervision mainly affects connectivity: removing it raises Dice and IoU by
0.08 and 0.27 points but lowers clDice by 0.26 points.  Fusion is optional:
removing it lowers Dice, IoU, and clDice by 0.12, 0.11, and 0.26 points but
improves HD95 by 0.36 pixels, so most of the gain over gating+fusion comes
from the fixed selector rather than fusion removal.  We keep the full model
because it is best on Dice, IoU, and clDice; the no-fusion variant is a
competitive alternative when boundary accuracy is prioritized.  No component
dominates every metric, but fixed deployment provides the largest single gain.
